\section{Séries temporais e redes recorrentes (LSTM)}
\label{sec:lstm}

\subsection{Por que uma MLP comum não basta}
Uma MLP trata cada entrada como um conjunto de números \emph{sem ordem}. Mas em
uma série temporal a \textbf{ordem é a informação}: ``subiu três dias e caiu no
quarto'' é diferente de ``caiu e subiu três''. Precisamos de uma rede que tenha
\emph{memória} da sequência.

\subsection{Redes recorrentes (RNN)}
\begin{definicao}[RNN]
Uma rede recorrente processa a sequência passo a passo, carregando um
\textbf{estado} (uma ``memória'') que é atualizado a cada novo elemento. Assim, a
saída no dia $t$ depende não só de $x_t$, mas de tudo que veio antes.
\end{definicao}

RNNs simples têm um problema: ao processar sequências longas, o sinal de
aprendizado ou \emph{some} (gradiente que desaparece) ou \emph{explode}, e a rede
``esquece'' o que viu há muitos passos.

\subsection{LSTM: memória com portões}
A \textbf{LSTM} (\emph{Long Short-Term Memory}) resolve isso com uma célula de
memória controlada por \textbf{portões} (\emph{gates}) --- pequenas
sub-redes que decidem, a cada passo:
\begin{itemize}[noitemsep]
  \item \textbf{esquecer}: o que descartar da memória antiga;
  \item \textbf{entrada}: o que da informação nova guardar;
  \item \textbf{saída}: o que da memória expor como resultado do passo.
\end{itemize}

\begin{intuicao}
Pense na LSTM como um leitor com um caderno de anotações. A cada frase ele
decide o que apagar, o que anotar de novo e o que usar agora. Esse controle
ativo da memória deixa a LSTM lembrar de padrões relevantes por muito mais tempo
que uma RNN simples --- ideal para capturar regimes de volatilidade que se
estendem por semanas.
\end{intuicao}

\subsection{Janelas deslizantes: fatiando o tempo}
A LSTM recebe \emph{sequências de tamanho fixo}. Convertemos a série contínua em
pedaços com a técnica da \textbf{janela deslizante}: uma janela de comprimento
$L$ varre a série de uma em uma posição.

\begin{exemplo}
Série $[a,b,c,d,e]$ com janelas de tamanho 3 vira:
$[a,b,c],\ [b,c,d],\ [c,d,e]$. Cada janela é uma ``frase'' de 3 dias que a LSTM
lê inteira.
\end{exemplo}

\begin{lstlisting}[caption={Gerando janelas de tamanho fixo a partir de um vetor 1D.}]
import numpy as np
def make_windows(valores, L=30, passo=1):
    idx = range(0, len(valores) - L + 1, passo)
    janelas = np.stack([valores[i:i+L] for i in idx])
    return janelas[..., np.newaxis]   # forma (n_janelas, L, 1)
\end{lstlisting}

A forma final \texttt{(n\_janelas, L, 1)} é o que a LSTM espera: $n$ exemplos, cada
um com $L$ passos no tempo e $1$ variável por passo (o log-retorno).

\begin{naprat}
Usamos janelas de \textbf{30 passos} (\texttt{window\_size=30}) --- cerca de seis
semanas úteis ---, valor consagrado em implementações de referência de detecção
de anomalias com LSTM. As janelas são geradas \emph{dentro} de cada partição
(treino e teste separadamente), para que nenhuma janela cruze a fronteira
temporal e reintroduza vazamento. Uma consequência (importante na avaliação): como
as janelas se sobrepõem, um único dia anômalo aparece em até 30 janelas
consecutivas --- uma anomalia vira uma ``rajada'' de marcações.
\end{naprat}
